\section{A canonical flow from the endpoint momenta}
\label{sec:v77-canonical}
The roof field can be specified on the whole regular exact event,
without choosing critical disks or endpoint caps. The particular
section in \eqref{eq:v76-event} is useful here: it bounds both endpoint
momenta away from grazing, even though no internal incidence or
clearance is bounded away from zero.

For one regular physical word use locally lifted endpoint angles
$q=(\alpha_0,\alpha_m)$, and write $F(q)$ for its reduced length.
The endpoint map is locally invertible, since the optical twist
$B_m=\partial_{p_0}\alpha_m$ is nonzero. Set
\[
 g=\nabla_qF=R(-p_0,p_m),\qquad K=D_q^2F,
 \qquad \lambda_\circ=\frac{9\sqrt{19}}{200},\quad
 \Lambda_\circ=\frac{47}{6}.
\]
The terminal tangential velocity has the same value before and after
reflection, which fixes the sign in the second component of $g$.
All differentiations in this section are on actual regular words.

\begin{lemma}[Uniform endpoint-angle Hessians on the exact event]
\label{lem:v77-hessian}
At every regular physical trajectory,
\begin{equation}\label{eq:v77-general-hessian}
 R\operatorname{diag}(c_0,c_m)\le K
       \le (2R_+^2/\tau_-+R_+)I_2.
\end{equation}
On the original exact event, in particular,
\begin{equation}\label{eq:v77-uniform-hessian}
                  \lambda_\circ I_2\le K\le\Lambda_\circ I_2.
\end{equation}
These are pointwise quadratic-form inequalities, uniform in $m,R,n,k$.
They assert neither convexity of a disconnected physical domain nor
uniform control of derivatives of $K$.
\end{lemma}
\begin{proof}
At a regular trajectory let $w_0$ be its outgoing first-flight
velocity and $w_m$ the negative of its incoming last-flight velocity.
The polygonal length, minimized over the prescribed intermediate
closed disks, is supported at this trajectory by
\[
 F(q_*)-w_0\cdot(Q_0(q_0)-Q_0(q_{0,*}))
        -w_m\cdot(Q_m(q_m)-Q_m(q_{m,*})).
\]
To see this directly, the internal gradient at contact $i$ is
$-2c_i n_i$. Its scalar product with a displacement into that disk
is nonnegative. The convex gradient inequality for the polygonal
length therefore gives the displayed support after minimizing the
internal positions. Equality holds at the physical trajectory.
The internal Hessian is positive definite, so the reduced roof is
smooth on its regular neighborhood. Since $Q_i''=-Rn_i$ and
$w_i\cdot n_i=c_i$, differentiating the touching support proves the
lower inequality in \eqref{eq:v77-general-hessian}.

For the upper inequality keep the internal contacts fixed while
varying the endpoints. The transverse second variation of a flight
is its squared transverse displacement divided by its length.
The two endpoint speeds have norm $R$. Their contribution is at
most $2R^2|dq|^2/\tau_-$; this bound includes the case $m=1$,
when both displacements occur on the same flight. The endpoint
curvature terms contribute at most $R|dq|^2$. The reduced Hessian
is the Schur complement of the internal Hessian and cannot exceed
this trial quadratic form. Taking the indicated extreme radii gives
the upper bound, with $2R_+^2/\tau_-+R_+=47/6$.

In the five rectangles of $Y_R^*$ one has $|p|<9/10$. For the only
moving momentum centers this follows from
\[
 \frac{2(53/100)}{1+(53/100)^2}+\frac1{200}<\frac9{10};
\]
the remaining centers and the added rectangle satisfy the bound
immediately. Both the initial and terminal states belong to this
same section. Thus $c_0,c_m>\sqrt{19}/10$, which proves
\eqref{eq:v77-uniform-hessian}.
\end{proof}

Away from $g=0$ define the circle-valued physical observable
\begin{equation}\label{eq:v77-angle}
 \vartheta(x)=\arg(-p_0(x),p_m(x))\in\R/(2\pi\Z).
\end{equation}
It is not the normal angle of an auxiliary endpoint cap. The
orientation of each disk is the fixed orientation already used
for the collision coordinates. Put $e=g/|g|$ and
\begin{equation}\label{eq:v77-chi}
 \chi=e^{\mathsf T}\operatorname{cof}(K)e,
 \qquad v=K^{-1}g,\qquad s=g^{\mathsf T}K^{-1}g.
\end{equation}
Here $\operatorname{cof}(K)$ is the adjugate matrix; $K$ is symmetric.
In dimension two its eigenvalues are those of $K$ in reverse order.
In particular $\lambda_\circ\le\chi\le\Lambda_\circ$ on the exact event.

\begin{theorem}[An explicit critical-safe forward roof field]
\label{thm:v77-field}
There is a uniquely specified smooth vector field $X$ on the regular
exact event with its interior critical points removed such that
\begin{equation}\label{eq:v77-field-defining}
                          XL=1,\qquad X\vartheta=0.
\end{equation}
In endpoint-angle coordinates it is
\begin{equation}\label{eq:v77-newton-field}
                              X_q=\frac{K^{-1}g}{s}.
\end{equation}
In any local lift of $\vartheta$ on the original collision cylinder,
it is equivalently
\[
 X=\frac{\vartheta_p\partial_\alpha-
                   \vartheta_\alpha\partial_p}
          {L_\alpha\vartheta_p-L_p\vartheta_\alpha}.
\]
No source density, guard, partition, or critical radius enters its
choice. Stop its flow at the first exit from the regular exact event.
For a forward survivor and $r_h=|g(\Phi_hx)|$, $r_0=|g(x)|$, one has
\begin{align}
 2\lambda_\circ h&\le r_h^2-r_0^2\le2\Lambda_\circ h,
                                      \label{eq:v77-radial-speed}\\
 \operatorname{length}(q([0,h]))
       &\le\frac{r_h-r_0}{\lambda_\circ}
        \le\frac{\sqrt{2\Lambda_\circ h}}{\lambda_\circ},
                                      \label{eq:v77-endpoint-motion}\\
 |p_i(\Phi_hx)-p_i(x)|&\le
       \frac{\sqrt{2\Lambda_\circ h}}{R_-},\qquad i=0,m.
                                      \label{eq:v77-momentum-motion}
\end{align}
The same estimates hold up to the limiting time of a finite lifetime.
A forward trajectory cannot terminate at a regular interior critical
point of $L$. The estimates do not assert survival through a true
singular, first-hit, or section-decision boundary.
\end{theorem}
\begin{proof}
Writing $g^\perp=(-g_2,g_1)$ gives
$d\vartheta=(g^\perp)^{\mathsf T}K\,dq/|g|^2$ and hence
\begin{equation}\label{eq:v77-jacobian}
       \det D_q(F,\vartheta)=\frac{g^{\mathsf T}
                   \operatorname{cof}(K)g}{|g|^2}=\chi>0.
\end{equation}
The endpoint map has Jacobian $B_m$, so
$|\det D_{(\alpha,p)}(L,\vartheta)|=|B_m|\chi>0$.
The two equations \eqref{eq:v77-field-defining} therefore specify a
unique field. Formula \eqref{eq:v77-newton-field} satisfies them:
$g\cdot X_q=1$ and $KX_q=g/s$ is radial in $g$.
Endpoint lifts differ by translations by multiples of $2\pi$;
local lifts of $\vartheta$ also differ by constants. Both formulas
therefore agree on overlaps and define the field without gluing.

Along the flow $\dot g=g/s$. Its direction is constant and
$\dot r=r/s>0$. Since
$r^2/\Lambda_\circ\le s\le r^2/\lambda_\circ$,
one has $2\lambda_\circ\le d(r^2)/dh\le2\Lambda_\circ$.
This proves \eqref{eq:v77-radial-speed}. Using $r$ as the parameter,
\[
                     \frac{dq}{dr}=K^{-1}e.
\]
Integration gives \eqref{eq:v77-endpoint-motion}, including length
rather than merely endpoint displacement. The direction of $g$ is
constant, so $|g(\Phi_hx)-g(x)|=r_h-r_0$ and the first-variation
formula gives \eqref{eq:v77-momentum-motion}. In particular $r_h$
cannot tend to zero in forward time. On a compact regular subset
away from $g=0$ the ordinary continuation theorem applies; any
remaining finite exit is a true boundary of the regular exact event.
\end{proof}

\begin{proposition}[The cancellation in the physical drift]
\label{prop:v77-drift}
Let $T=D^3F$ and let $a$ be a locally Lipschitz density relative to
$dq$. Away from $g=0$, in the almost-everywhere sense,
\begin{equation}\label{eq:v77-drift-formula}
 \operatorname{div}_q(aX_q)
 =\frac{\nabla a\cdot v}{s}
   +a\left\{\frac{T[v,v,v]}{s^2}
              -\frac{v\cdot\nabla\log\det K}{s}\right\}.
\end{equation}
On a regular neighborhood where $\|T\|\le T_0$, this gives
\begin{equation}\label{eq:v77-drift-local-bound}
 |\operatorname{div}_q(aX_q)|
 \le\frac1{|g|}\left\{
 \frac{\Lambda_\circ}{\lambda_\circ}|\nabla a|
 +\left(\frac{\Lambda_\circ^2}{\lambda_\circ^3}
       +\frac{2\Lambda_\circ}{\lambda_\circ^2}\right)|a|T_0\right\}.
\end{equation}
In particular the constant-density quadratic part has zero divergence.
The bounds on $T$ and $\nabla a$ in this assertion are local, not
collision-uniform conclusions.
\end{proposition}
\begin{proof}
Differentiation of $v=K^{-1}g$ yields
$\operatorname{div}v=2-v\cdot\nabla\log\det K$.
Also $v\cdot\nabla s=2s-T[v,v,v]$. The two terms $2/s$
cancel upon differentiating $a v/s$, proving
\eqref{eq:v77-drift-formula}. Use $|v|\le|g|/\lambda_\circ$,
$s\ge|g|^2/\Lambda_\circ$, and
$|\nabla\log\det K|\le2T_0/\lambda_\circ$ for the bound.
Near a regular interior critical point the uniform positive Hessian
makes $|g|$ comparable to distance, so the remaining displayed
singularity is locally integrable in two dimensions. This last
observation concerns a local integral only.
\end{proof}
